\documentclass[11pt]{article}
\newcommand{\paso}{2}
\usepackage[T1]{fontenc}
\usepackage[utf8]{inputenc}
\usepackage[spanish,es-nodecimaldot,es-noquoting,es-lcroman]{babel}
\spanishplainpercent
\usepackage{lmodern}
\usepackage[a4paper,margin=21mm,headheight=15pt]{geometry}
\usepackage{amsmath,amssymb,bm}
\usepackage{booktabs,array}
\usepackage{tikz}
\usetikzlibrary{arrows.meta,positioning,calc}
\usepackage[most]{tcolorbox}
\usepackage{microtype}
\usepackage{fancyhdr}
\usepackage{hyperref}
\hypersetup{colorlinks=true,linkcolor=blue!45!black,urlcolor=blue!45!black}
\setlength{\parindent}{0pt}
\setlength{\parskip}{5pt}
\setlength{\emergencystretch}{2em}
\pagestyle{fancy}
\fancyhf{}
\fancyhead[L]{\small BEM del cilindro: paso \paso}
\fancyhead[R]{\small Guía matemática independiente}
\fancyfoot[C]{\thepage}
\newcommand{\vct}[1]{\bm{#1}}
\newcommand{\dd}{\,\mathrm{d}}
\newcommand{\ii}{\mathrm{i}}
\newcommand{\norm}[1]{\left\lVert#1\right\rVert}
\newcommand{\VP}{\operatorname{VP}}
\newcommand{\tr}{\operatorname{tr}}
\newtcolorbox{idea}[1]{breakable,colback=blue!3,colframe=blue!50!black,
  title={#1},fonttitle=\bfseries,boxrule=0.7pt}
\newtcolorbox{alerta}[1]{breakable,colback=orange!4,colframe=orange!65!black,
  title={#1},fonttitle=\bfseries,boxrule=0.7pt}
\newtcolorbox{resultado}[1]{breakable,colback=teal!3,colframe=teal!55!black,
  title={#1},fonttitle=\bfseries,boxrule=0.7pt}
\newcommand{\inicio}[2]{%
 {\small\color{blue!55!black}\bfseries SERIE: LOS SIETE PASOS MATEMÁTICOS DEL BEM}\par
 {\LARGE\bfseries Paso \paso. #1\par}\vspace{3mm}
 \begin{idea}{¿Qué busca este paso?}#2\end{idea}}
\newcommand{\mapa}[3]{%
 \begin{center}
 \begin{tikzpicture}[node distance=4mm,
   caja/.style={draw=blue!50!black,fill=blue!3,rounded corners=2pt,
     text width=13cm,align=center,inner sep=6pt},
   >={Stealth[length=2mm]}]
 \node[caja] (a){#1};
 \node[caja,below=of a] (b){#2};
 \node[caja,below=of b] (c){#3};
 \draw[->,thick,blue!60!black] (a)--(b);
 \draw[->,thick,blue!60!black] (b)--(c);
 \end{tikzpicture}
 \end{center}}
\newcommand{\base}{%
 \smallskip{\footnotesize\color{black!65}
 Modelo de referencia: cilindro sólido, material homogéneo e isótropo,
 deformación plana, pequeñas deformaciones, sin fuerzas de volumen;
 amplitudes complejas con \(e^{+\ii\omega t}\), \(\omega>0\).}\par}
\newcommand{\fuentes}{%
 \par\medskip\begingroup\footnotesize
 \textbf{Referencia y continuación.}
 Este documento desarrolla una etapa del esquema
 \texttt{esquema\_ecuaciones\_bem.tex}. Las ecuaciones y los ejemplos se
 contrastaron con \texttt{bem\_cilindro.py}, ambos en la carpeta superior.
 La guía general y sus deducciones ampliadas se encuentran en
 \texttt{../guia\_bem/guia\_bem.pdf}. No se modifica la implementación.
 \par Para recompilar este documento, mantener \texttt{estilo.tex} en
 la misma carpeta y ejecutar \texttt{pdflatex} sobre su archivo
 \texttt{.tex}, desde esa carpeta.\par\endgroup}

\begin{document}
\inicio{El material complejo}{Convertir las propiedades del material en los
coeficientes que necesita la elastodinámica: módulos de Lamé complejos y
velocidades de ondas P y S. El objetivo es representar rigidez y disipación
sin añadir nuevos grados de libertad al problema de frontera.}
\base
\mapa{Datos: \(\rho,\mu_0,\nu,\beta\)}
{Constitución: \(\mu^*=\mu_0(1+2\ii\beta)\), \(\lambda^*=2\mu^*\nu/(1-2\nu)\)}
{Propagación: \(c_s^2=\mu^*/\rho\), \(c_p^2=(\lambda^*+2\mu^*)/\rho\)}

\section*{1. Qué significa cada dato}
\begin{center}
\begin{tabular}{p{0.17\linewidth}p{0.70\linewidth}}
\toprule
Dato & Interpretación\\
\midrule
\(\rho>0\) & Masa por volumen. Controla la resistencia a acelerar.\\
\(\mu_0>0\) & Módulo de corte de referencia [Pa].\\
\(\nu\) & Coeficiente de Poisson real. Para el material estable usual:
\(-1<\nu<1/2\).\\
\(\beta\ge0\) & Parámetro adimensional de amortiguamiento del modelo.\\
\bottomrule
\end{tabular}
\end{center}
Los valores de referencia son
\[
 \rho=100\ \mathrm{kg/m^3},\quad \mu_0=10^6\ \mathrm{Pa},\quad
 \nu=0.25,\quad\beta=0.05.
\]
\begin{idea}{Rigidez y disipación cumplen funciones distintas}
La parte real del módulo relaciona la respuesta elástica con la deformación.
La parte imaginaria representa el desfase disipativo en el régimen armónico.
No significa que haya dos materiales distintos.
\end{idea}

\section*{2. Entrada y salida}
Este paso recibe constantes físicas reales y devuelve parámetros que se
reutilizan en cada frecuencia. En este modelo \(\mu^*\) y \(\lambda^*\) no
dependen explícitamente de \(\omega\). La dependencia en frecuencia entra en
los números de onda y en los núcleos del paso 4.

\newpage
\section*{3. Por qué un módulo complejo representa desfase}
Para verlo sin notación tensorial, pensemos en un corte simple con
deformación ingenieril \(\gamma(t)=\gamma_0\cos\omega t\). Su amplitud es real
y el esfuerzo es
\[
 \tau(t)=\operatorname{Re}\{\mu^*\gamma_0e^{\ii\omega t}\}.
\]
Si \(\mu^*=\mu'+\ii\mu''\), al multiplicar:
\[
 \tau(t)=\mu'\gamma_0\cos\omega t-\mu''\gamma_0\sin\omega t.
\]
El término con seno no está en fase con la deformación. Integrando la potencia
por unidad de volumen en un período \(T=2\pi/\omega\):
\[
 \Delta W=\int_0^T\tau(t)\dot\gamma(t)\dd t
          =\pi\mu''\gamma_0^2.
\]
Con la convención \(e^{+\ii\omega t}\), \(\mu''>0\) da disipación positiva.

\begin{resultado}{Cómo entra el parámetro del notebook}
\[
 \mu^*=\mu_0(1+2\ii\beta),\qquad
 \mu'=\mu_0,\quad\mu''=2\beta\mu_0,\quad
 \eta_{\text{pérdida}}=\frac{\mu''}{\mu'}=2\beta.
\]
Para \(\beta=0.05\), el factor de pérdida es \(0.10\).
La identificación de \(\beta\) con una razón de amortiguamiento modal es
aproximada y depende del contexto; no son definiciones universales idénticas.
\end{resultado}

\section*{4. De \(\mu^*,\nu\) a \(\lambda^*\)}
En elasticidad isótropa,
\[
 \nu=\frac{\lambda^*}{2(\lambda^*+\mu^*)}.
\]
Despejando:
\[
 2\nu\lambda^*+2\nu\mu^*=\lambda^*
 \quad\Longrightarrow\quad
 \boxed{\lambda^*=\frac{2\nu\mu^*}{1-2\nu}.}
\]
Como se conserva \(\nu\) real, ambos módulos comparten el mismo factor
\(1+2\ii\beta\). Para \(\nu=0.25\), la fracción vale uno y
\(\lambda^*=\mu^*\).

\begin{alerta}{Cerca de la incompresibilidad}
Si \(\nu\to1/2\), \(\lambda^*\) crece mucho. No se debe evaluar la fórmula en
\(\nu=1/2\) ni confundir una división por cero con un detalle numérico.
El límite incompresible requiere una formulación adecuada.
\end{alerta}

\newpage
\section*{5. De dónde salen las dos velocidades}
Tomemos una onda plana con vector de onda \(k\vct d\), donde \(\vct d\) es
unitario:
\[
 \vct u=\vct a\,e^{\ii\omega t-\ii k\vct d\cdot\vct x}.
\]
Al sustituir en Navier:
\[
 (\rho\omega^2-\mu^*k^2)\vct a
 -(\lambda^*+\mu^*)k^2\vct d(\vct d\cdot\vct a)=\vct0.
\]
\subsection*{5.1 Polarización transversal: onda S}
Si \(\vct d\cdot\vct a=0\), se anula el segundo término:
\[
 \rho\omega^2=\mu^*k_s^2,\qquad
 \boxed{c_s=\frac{\omega}{k_s}=\sqrt{\frac{\mu^*}{\rho}}.}
\]
\subsection*{5.2 Polarización longitudinal: onda P}
Si \(\vct a\) es paralela a \(\vct d\), los términos elásticos se suman:
\[
 \rho\omega^2=(\lambda^*+2\mu^*)k_p^2,\qquad
 \boxed{c_p=\sqrt{\frac{\lambda^*+2\mu^*}{\rho}}.}
\]
Las ondas S describen distorsión transversal; las P incluyen compresión
longitudinal. Las dos intervienen en una solución fundamental elástica,
aunque la respuesta global bajo esta carga sea radial.

\section*{6. Raíces complejas y atenuación}
Se usan raíces compatibles con la convención temporal y el campo saliente.
En el caso de referencia, la raíz principal tiene partes real e imaginaria
positivas. Por ello \(k=\omega/c\) tiene parte imaginaria negativa:
\[
 e^{-\ii kr}=e^{-\ii\operatorname{Re}(k)r}
              e^{\operatorname{Im}(k)r}.
\]
La amplitud disminuye con \(r\). Elegir arbitrariamente el signo opuesto
puede producir una onda entrante o creciente.

\begin{alerta}{Una velocidad compleja no es una rapidez medible única}
La rapidez de fase se asocia a \(\omega/\operatorname{Re}(k)\); la atenuación
a \(-\operatorname{Im}(k)\). No se deben interpretar directamente
\(\operatorname{Re}(c)\) y \(\operatorname{Im}(c)\) como dos velocidades físicas.
\end{alerta}

\newpage
\section*{7. Ejemplo numérico completo}
\[
 \mu^*=10^6+\ii10^5\ \mathrm{Pa},\qquad
 \lambda^*=10^6+\ii10^5\ \mathrm{Pa}.
\]
\[
 \begin{aligned}
 c_s&\simeq100.124611+\ii4.993777\ \mathrm{m/s},\\
 c_p&\simeq173.420914+\ii8.649476\ \mathrm{m/s}.
 \end{aligned}
\]
Una relación útil para los núcleos es
\[
 \gamma^2=\frac{c_s^2}{c_p^2}
 =\frac{\mu^*}{\lambda^*+2\mu^*}
 =\frac{1-2\nu}{2(1-\nu)}.
\]
Aquí \(\gamma^2=1/3\), pese a que las dos velocidades sean complejas.
Esto ocurre porque los módulos tienen el mismo factor complejo.

\section*{8. Controles y límites del procedimiento}
\begin{itemize}
 \item Con \(\beta=0\), los módulos y las velocidades vuelven a ser reales
 positivos para el material estable considerado.
 \item Comprobar \(\rho c_s^2=\mu^*\) y
 \(\rho c_p^2=\lambda^*+2\mu^*\); no comparar solo partes reales.
 \item Las unidades de \(c^2\) son \(\mathrm{m^2/s^2}\), pues
 \([\mu/\rho]=\mathrm{Pa}/(\mathrm{kg/m^3})\).
 \item No introducir \(\beta=5\) para indicar cinco por ciento:
 el valor adimensional es \(0.05\).
 \item Este es un modelo histerético de frecuencia positiva. Su extensión a
 todo el eje de frecuencias y a una respuesta transitoria causal no se obtiene
 usando sin más el mismo módulo complejo constante.
 \item Cambiar a \(e^{-\ii\omega t}\) exige cambiar coherentemente los signos
 imaginarios y la selección de ondas; no solo un exponente aislado.
\end{itemize}

\begin{resultado}{Qué se conserva para el siguiente paso}
\(\rho,\nu,\mu^*,\lambda^*,c_s,c_p\) describen el medio auxiliar homogéneo.
El paso 3 usa su reciprocidad; el paso 4 construye sus soluciones
fundamentales con \(z_s=\ii\omega r/c_s\) y \(z_p=\ii\omega r/c_p\).
\end{resultado}
\fuentes
\end{document}
